\documentclass{article}
\usepackage{geometry}
\usepackage{graphicx, url, hyperref, amsmath, amssymb, mathtools, comment, xcolor}
\usepackage{amsthm}
\usepackage[most]{tcolorbox}
\tcbuselibrary{skins, breakable}
\usepackage[shortlabels]{enumitem}
\usepackage{mathscinet} % for \Dbar
\usepackage{caption}
\usepackage{booktabs}
\usepackage{tabularx}

% Geometry settings
\geometry{margin=1.5in}

% --- Theorems and Styles ---
\theoremstyle{plain}
\newtheorem{theorem}{Theorem}
\newtheorem{lemma}{Lemma}
\newtheorem{claim}{Claim}

\theoremstyle{definition}
\newtheorem*{example}{Example}
\newtheorem{remark}{Remark}[section]
\newtheorem{definition}[theorem]{Definition}

% --- Custom Commands ---
\newcommand{\Aletheia}{\emph{Aletheia}}
\renewcommand{\comment}[1]{\marginpar{{\tiny{#1}\normalfont\par}}}
\newcommand{\tony}[1]{{\color{red}Tony: #1}}

% --- Defined Colors ---
\definecolor{boxblue}{RGB}{0, 0, 150}
\definecolor{boxback}{RGB}{245, 245, 255}

% --- The Problem Environment ---
\newtcolorbox{problem}[1]{%
    colback=boxback,
    colframe=boxblue,
    fonttitle=\bfseries\large,
    title={#1},
    sharp corners,
    enhanced,
    attach boxed title to top left={yshift=-2mm, xshift=2mm},
    boxed title style={colframe=boxblue, colback=boxblue},
    before skip=15pt plus 2pt,
    after skip=15pt plus 2pt,
    top=10pt, bottom=10pt, left=10pt, right=10pt
}

\newtcolorbox{solution}[1]{%
    colback=white,
    colframe=boxblue,
    fonttitle=\bfseries\large,
    title={#1},
    sharp corners,
    enhanced jigsaw, % Better frame handling for page breaks than just 'enhanced'
    breakable,       % <--- Allows the box to split across pages
    attach boxed title to top left={yshift=-2mm, xshift=2mm},
    boxed title style={colframe=boxblue, colback=boxblue},
    before skip=15pt plus 2pt,
    after skip=15pt plus 2pt,
    top=10pt, bottom=10pt, left=10pt, right=10pt
}



% Define custom commands for readability\newcommand{\BB}{\mathbb{B}}
\newcommand{\BB}{\mathbb{B}}
\newcommand{\VV}{\mathbb{V}}
\newcommand{\Q}{\mathbb{Q}}
\newcommand{\Gr}{\mathrm{Gr}}
\newcommand{\aug}{\mathrm{aug}}
\newcommand{\GL}{\mathrm{GL}}
\newcommand{\frR}{\mathfrak{R}}
\DeclareMathOperator{\htop}{ht}

\DeclareMathOperator{\sgn}{sgn}
\DeclareMathOperator{\id}{id}
\newcommand{\PSp}{\mathrm{PSp}}
\newcommand{\Pf}{\mathrm{Pf}}
\DeclareMathOperator{\PSO}{PSO}

\date{}


\begin{document}
\emph{The prompts and responses below have been edited only for
formatting compatibility.}

\begin{problem}{User Prompt} We work over $\mathbb{C}$. Let $\mathbb{E} \to \mathcal{M}_{g}$ be the rank $g$ Hodge bundle over the moduli stack of genus $g$ curves. Prove that $\mathbb{E}$ does not contain any non-trivial sub-vector bundles.  (A sub-vector bundle is a vector bundle $V \to \mathcal{M}_{g}$ with an injective map $i : V \hookrightarrow \mathbb{E}$ which is injective on all fibers. A sub-vector bundle is trivial if it is either the 0-section or the whole bundle.)
\end{problem}

\begin{solution}{Model Response (\Aletheia)}
\begin{abstract}
We prove that the rank $g$ Hodge bundle $\mathbb{E} \to \mathcal{M}_g$ over the moduli stack of smooth complex algebraic curves of genus $g \ge 2$ admits no non-trivial proper sub-vector bundles. This is achieved by analyzing the equivariant characters of automorphisms on putative subbundles over specifically constructed loci of curves with a $(\mathbb{Z}/2\mathbb{Z})^2$-action. By extracting a rigid functional equation for the character traces, we categorically force the dimension of the isotypic invariant subspaces to be non-integral, yielding an algebraic contradiction.
\end{abstract}

\section{Introduction}
Let $\mathcal{M}_g$ be the moduli stack of smooth projective curves of genus $g \ge 2$ over $\mathbb{C}$. The Hodge bundle $\mathbb{E} \to \mathcal{M}_g$ is the rank $g$ vector bundle whose fiber over a moduli point $[C]$ is the space of holomorphic abelian differentials $H^0(C, \omega_C)$.

Suppose for the sake of contradiction that there exists a globally defined proper non-trivial sub-vector bundle $V \subset \mathbb{E}$ of rank $r$. By definition, this means $V \to \mathcal{M}_g$ admits a bundle map $V \hookrightarrow \mathbb{E}$ that is injective on all fibers, and its rank satisfies $0 < r < g$. For any curve $C \in \mathcal{M}_g$, the fiber $V_C$ is an $r$-dimensional subspace of $H^0(C, \omega_C)$ which is naturally invariant under the pullback action of the full automorphism group $\operatorname{Aut}(C)$.

\section{Equivariant Traces of Involutions}

\begin{lemma} \label{lem:constant_trace}
Let $\tau$ be an orientation-preserving involution on a smooth curve $C$ of genus $g$ with $2m$ fixed points. The trace of the pullback $\tau^*$ on $V_C$ is a constant integer $f(m)$ depending only on $m$ and the bundle $V$.
\end{lemma}

\begin{proof}
By the topological classification of surfaces, two orientation-preserving involutions on a closed orientable surface $\Sigma_g$ with the same number of fixed points correspond to double covers over a quotient surface of a uniquely determined genus. Thus, such involutions are conjugate in the mapping class group $\operatorname{Mod}_g$. The putative bundle $V$ functorially corresponds to a $\operatorname{Mod}_g$-equivariant holomorphic vector bundle over the Teichm\"uller space $\mathcal{T}_g$.

By the Nielsen Realization Theorem (S. P. Kerckhoff, \textit{The Nielsen Realization Problem}, Ann. of Math., 1983), the fixed-point locus $\mathcal{T}_g^\tau$ of $\tau$ is non-empty and naturally identifies with the Teichm\"uller space of the quotient orbifold $\Sigma_g / \langle \tau \rangle$. Since the Teichm\"uller space of any orbifold is homeomorphic to a Euclidean space, $\mathcal{T}_g^\tau$ is a connected manifold. The restriction of the equivariant bundle $V$ to $\mathcal{T}_g^\tau$ globally splits into $+1$ and $-1$ eigenbundles for $\tau^*$. Because $\mathcal{T}_g^\tau$ is connected, the ranks of these eigenbundles are strictly constant. The trace of $\tau^*$ on $V_X$ for any $X \in \mathcal{T}_g^\tau$ is exactly the difference of these constant ranks, forming an integer $f(m)$ dictated entirely by the conjugacy class of $\tau$, and hence uniquely determined by $m$.
\end{proof}

By the Riemann-Hurwitz formula applied to the quotient map $C \to C/\langle\tau\rangle$, we have $2g - 2 = 2(2g' - 2) + 2m$, meaning $m = g + 1 - 2g'$. Since the quotient genus $g'$ is a non-negative integer, the geometrically valid parameters $m$ lie strictly in the finite set:
\[ S = \{ m \in \mathbb{Z} \mid 0 \le m \le g+1, \; m \equiv g+1 \pmod 2 \}. \]

\section{A Functional Equation via $(\mathbb{Z}/2\mathbb{Z})^2$-Covers}

\begin{lemma} \label{lem:func_eq}
For any three integers $m_1, m_2, m_3 \in S$ such that $m_1 + m_2 + m_3 = g + 3$ and at most one $m_i = 0$, we have the uniform trace relation $f(m_1) + f(m_2) + f(m_3) = -r$.
\end{lemma}

\begin{proof}
Given a valid triple $(m_1, m_2, m_3)$, we choose generic, mutually coprime, square-free polynomials $A(x), B(x), C(x)$ over $\mathbb{C}$ of degrees $m_1, m_2, m_3$, respectively. Let $K = \mathbb{C}(x)$. We define a smooth projective curve $C$ as the smooth model of the function field $L = K(y, z)$ where $y^2 = A(x)C(x)$ and $z^2 = B(x)C(x)$. Because at most one $m_i=0$, the polynomials cannot all be constants; thus the extension $L/K$ forms a fully generated degree $4$ Galois extension with Galois group $G \cong (\mathbb{Z}/2\mathbb{Z})^2$, guaranteeing $C$ is a connected curve. Since $m_i \equiv g+1 \pmod 2$, the degrees of $A, B,$ and $C$ strictly share the same parity. Consequently, the polynomials $A(x)C(x)$ and $B(x)C(x)$ possess even degrees, meaning $L/K$ is unramified over the place at $x = \infty$.

By local Galois analysis, the ramification of $L/K$ is localized entirely over the roots of $A, B,$ and $C$:
\begin{itemize}
    \item Over a root of $A(x)$, the subextension $K(z)/K$ is unramified (as $B$ and $C$ do not vanish there), naturally splitting into $2$ places in $K(z)$. Adjoining $y$ introduces ramification of index $2$ at both places. Thus, there are $2$ places in $L$, both with ramification index $2$ over $K$.
    \item Symmetrically, there are exactly $2$ places over each root of $B(x)$, each with ramification index $2$ over $K$.
    \item Over a root $c$ of $C(x)$, the subextension $K(y/z)/K$ evaluates locally to $(y/z)^2 = A(c)/B(c) \neq 0$, which cleanly splits into $2$ places over $\mathbb{C}$. Adjoining $y$ (and therefore $z$) introduces ramification of index $2$ at both places.
\end{itemize}
By the Riemann-Hurwitz formula applied to $L/K$, the genus $g_C$ satisfies:
\[ 2g_C - 2 = 4(-2) + 2m_1 + 2m_2 + 2m_3 = -8 + 2(g+3) = 2g - 2, \]
which formally confirms $g_C = g$.

The curve $C$ admits an automorphism group $G \cong (\mathbb{Z}/2\mathbb{Z})^2$ generated by the involutions $\tau_1(y, z) = (-y, z)$ and $\tau_2(y, z) = (y, -z)$, with $\tau_3 = \tau_1 \tau_2 = (-y, -z)$. The fixed points of $\tau_1$ correspond precisely to the places of $L$ fixed by the involution $\tau_1$, meaning $\tau_1$ lies in their inertia group. Because $L = K(z)(y)$ and $y^2 = A(x)C(x)$, ramification in $L/K(z)$ occurs exactly where the local valuation of $A(x)C(x)$ in $K(z)$ is odd.

Over a root of $A(x)$, $A(x)C(x)$ has a simple zero in $K$, and since $K(z)/K$ is unramified here, the valuation in $K(z)$ remains $1$ (odd). Thus, these places strictly ramify in $L/K(z)$, meaning $\tau_1$ generates the inertia group and fixes the places. As there are $2$ places in $K(z)$ over each of the $m_1$ roots of $A(x)$, this contributes exactly $2m_1$ fixed points.

Conversely, at a place of $K(z)$ lying over a root $c$ of $C(x)$, the local affine parameter for the base field $K = \mathbb{C}(x)$ is $x-c$. We have $z^2 = B(x)C(x)$. Since $B(c) \neq 0$ and $C(x)$ has a simple root at $c$, $z^2$ vanishes to order $1$ with respect to $x-c$, making $z$ a local uniformizer for $K(z)$ at this unique place. Consequently, $x-c$ vanishes to order $2$ with respect to $z$, and thus $A(x)C(x)$ also inherently vanishes to order $2$. Since its valuation is even, $(y/z)^2$ evaluates locally to a non-zero complex number $A(c)/B(c)$. Extracting the square root cleanly splits into two unramified places in $L$. Because they are unramified, the non-trivial automorphism $\tau_1$ of $L/K(z)$ simply swaps these two places, yielding zero fixed points.

Finally, over a root of $B(x)$, the field $K(y)$ is unramified, meaning there are $2$ places in $K(y)$ where $y$ takes distinct non-zero values. The involution $\tau_1$ maps $y \mapsto -y$, explicitly swapping these places, thus fixing none.

Therefore, $\tau_1$ possesses exactly $2m_1$ fixed points. By identical symmetry, $\tau_2$ and $\tau_3$ have exactly $2m_2$ and $2m_3$ fixed points, respectively.

The space $V_C$ decomposes into four $G$-isotypic components: $V^{++}, V^{+-}, V^{-+}, V^{--}$. Projecting onto the invariant subspace via character theory yields:
\[ \dim V^{++} = \frac{1}{4} \big( \operatorname{tr}(1) + \operatorname{tr}(\tau_1) + \operatorname{tr}(\tau_2) + \operatorname{tr}(\tau_3) \big) = \frac{1}{4} \big( r + f(m_1) + f(m_2) + f(m_3) \big). \]
Crucially, $V^{++}$ is a subspace of the $G$-invariant differentials $H^0(C, \omega_C)^G$, which canonically correspond to holomorphic differentials on the quotient curve $C/G$. Since the invariant field is exactly $K = \mathbb{C}(x)$, the quotient curve is natively $\mathbb{P}^1$. Because $H^0(\mathbb{P}^1, \omega_{\mathbb{P}^1}) = 0$, we unconditionally deduce $\dim V^{++} = 0$. This forces the functional equation $f(m_1) + f(m_2) + f(m_3) = -r$.
\end{proof}

\section{Solving the Character Recurrence}

\begin{lemma} \label{lem:recurrence_solution}
For all $m \in S$, the unique solution to the functional equation is $f(m) = \frac{r(1-m)}{g}$.
\end{lemma}

\begin{proof}
Let $y(m) = f(m) - \frac{r(1-m)}{g}$. Since the affine expression $\frac{r(1-m)}{g}$ algebraically satisfies the sum relation for any triple summing to $g+3$, we immediately deduce the homogeneous equation $y(m_1) + y(m_2) + y(m_3) = 0$ for all valid triples.

When $m = g+1$, the involution possesses $2g+2$ fixed points, uniquely identifying it as the hyperelliptic involution. The hyperelliptic involution universally acts as $-I$ on the canonical bundle, thus acting as $-I$ on $H^0(C, \omega_C)$ and its subspace $V_C$. Hence, $f(g+1) = -r$, enforcing the strong boundary condition $y(g+1) = 0$.

\textbf{Case $g$ is even:} The set $S$ consists of odd integers $m_i = 2k_i+1$ with $0 \le k_i \le g/2 = n$. Since $m_i \ge 1$, no valid index is zero. Let $y_k = y(2k+1)$. The relation for the index triple $(k_1, k_2, k_3) = (0, n, 0)$ corresponds to the valid $m$-values $(1, g+1, 1)$, which cleanly sum to $g+3$. This yields $2y_0 + y_n = 0 \implies y_0 = 0$. For any $0 \le j \le n-1$, the index triples $(1, j, n-1-j)$ and $(0, j+1, n-1-j)$ correspond to valid $m$-triples summing to $g+3$. Subtracting their respective sum equations isolates $y_{j+1} - y_j = y_1$, forming an arithmetic progression $y_k = k y_1$. Because $y_n = 0$ and $n \ge 1$, we firmly deduce $y_k = 0$ for all $k$.

\textbf{Case $g$ is odd:} The set $S$ consists of even integers $m_i = 2k_i$ with $0 \le k_i \le (g+1)/2 = n-1$, where $n = (g+3)/2$. Let $y_k = y(2k)$. The boundary condition guarantees $y_{n-1} = 0$. For $1 \le k \le n-2$, the index triples $(1, k, n-1-k)$ and $(0, k+1, n-1-k)$ correspond to $m$-triples summing to $g+3$, each structurally containing at most one zero index. Subtracting their respective equations gives $y_{k+1} - y_k = y_1 - y_0$, establishing an arithmetic progression. The valid index triple $(0, 1, n-1)$ corresponds to $m$-values $(0, 2, g+1)$, safely possessing exactly one zero, and yields $y_0 + y_1 + y_{n-1} = 0 \implies y_1 = -y_0$. The progression thus evaluates to $y_k = y_0(1-2k)$. At $k=n-1$, $y_{n-1} = y_0(1-2(n-1)) = y_0(3-2n) = 0$. Since $g \ge 3$, we have $3-2n = -g \neq 0$, forcing $y_0 = 0$, meaning $y_k = 0$ for all $k$.
\end{proof}

\section{The Integrality Contradiction}

\begin{theorem} \label{thm:no_subbundles}
The Hodge bundle $\mathbb{E} \to \mathcal{M}_g$ contains no non-trivial proper sub-vector bundles.
\end{theorem}

\begin{proof}
Isolating the trace projection for the $V^{+-}$ eigenspace yields:
\[ \dim V^{+-} = \frac{1}{4} \big( \operatorname{tr}(1) + \operatorname{tr}(\tau_1) - \operatorname{tr}(\tau_2) - \operatorname{tr}(\tau_3) \big) = \frac{1}{4} \big( r + f(m_1) - f(m_2) - f(m_3) \big). \]
Substituting the algebraic identity $-f(m_2) - f(m_3) = r + f(m_1)$ via Lemma \ref{lem:func_eq} cleanly simplifies this structure to:
\[ \dim V^{+-} = \frac{1}{2} \big( r + f(m_1) \big). \]
Consequently, $2 \dim V^{+-} = r + f(m_1) = r + \frac{r(1-m_1)}{g}$.
Because $\dim V^{+-}$ represents the formal dimension of a complex vector space, $2 \dim V^{+-}$ must intrinsically be an even integer. Thus, the evaluated expression $r + \frac{r(1-m_1)}{g}$ is restricted to be an \textbf{even integer} for all values $m_1 \in S$ capable of completing a valid triple.

\begin{itemize}
    \item \textbf{If $g$ is even:} Both $1$ and $3$ inherently belong to $S$. They seamlessly extend to the valid triples $(1, 1, g+1)$ and $(3, 1, g-1)$ respectively, summing to $g+3$ with no zero entries. Evaluating at $m_1=1$ strictly produces $r$, logically demanding that $r$ is an even integer. Evaluating at $m_1=3$ yields $r - \frac{2r}{g}$, which must simultaneously act as an even integer.

    \item \textbf{If $g$ is odd:} Both $0$ and $2$ inherently belong to $S$. They smoothly extend to the valid triples $(0, 2, g+1)$ and $(2, 2, g-1)$ respectively, summing to $g+3$ with at most one zero entry (since $g \ge 3 \implies g-1 \ge 2$). Evaluating at $m_1=0$ yields $r + \frac{r}{g}$, which is forced to be an even integer. Evaluating at $m_1=2$ yields $r - \frac{r}{g}$, which must similarly act as an even integer.
\end{itemize}

In either topological scenario, calculating the arithmetic difference between the evaluations at the two adjacent valid states strictly isolates $\frac{2r}{g}$. Because the difference between any two even integers categorically generates another even integer, we deduce:
\[ \frac{2r}{g} \in 2\mathbb{Z} \implies \frac{r}{g} \in \mathbb{Z} \implies r \equiv 0 \pmod g. \]
Since $V$ was hypothesized to be a proper subbundle, its rank is structurally constrained to the open interval $0 < r < g$. However, the parity condition categorically demands that $r$ behaves as a positive integer multiple of $g$. As absolutely no multiple of $g$ exists strictly within the open interval $(0, g)$, we achieve an absolute mathematical contradiction.

Hence, no non-trivial proper sub-vector bundle $V \subset \mathbb{E}$ can topologically or algebraically exist over $\mathcal{M}_g$.
\end{proof}
\end{solution}


\end{document}
